\section{Изотропия, обратимость, симметрии}

\begin{axiom}[A12 · Изотропия / Lorentz на $\Tlayer$]
\label{ax:A12}
Микрошаг изотропен: веса на соседях $N(x)$ равны.
Макроскопический свет на~$\Tlayer$ изотропен.
\end{axiom}

\begin{axiom}[A13 · Обратимость микрошага]
\label{ax:A13}
$g$ обратим на дискретном алфавите: $\exists g^{-1}$, $g^{-1}(g(\symCapitalPsi))=\symCapitalPsi$.

\end{axiom}

\begin{axiom}[A14 · Дискретные симметрии P, C, T, U(1)$_{\mathrm{vac}}$]
\label{ax:A14}
На конфигурациях $\symCapitalPsi=\{z(x)\}_{x\in\carrierLattice}$ и истории «лягушки'' $(Z_t,Z_{t-1})$
заданы преобразования:
\begin{itemize}[nosep]
\item \textbf{$U(1)_{\mathrm{vac}}$:} $z(x)\mapsto e^{i\theta}z(x)$ для всех $x\in\carrierLattice$
      ($\theta$ --- одна глобальная фаза);
\item \textbf{$P$ (parity):} пространственная инверсия $x\mapsto Px$ на~$\carrierLattice$,
      $(P\symCapitalPsi)(x):=z(Px)$;
\item \textbf{$C$ (charge conjugation):} комплексное сопряжение $z\mapsto z^*$
      на~$\mathbb{Z}_{N_{\mathrm{ring}}}[i]$;
\item \textbf{$T$ (time reversal):} обращение хода «лягушки'' и сопряжение истории,
      $(TZ_t,TZ_{t-1}):=(Z_{t-1}^*,Z_t^*)$.
\end{itemize}
Локальный закон~$g$ \emph{согласован} с ними:
\begin{align}
  g(U(1)_\theta\symCapitalPsi) &= U(1)_\theta\,g(\symCapitalPsi), \label{eq:A14-U1}\\
  g(P\symCapitalPsi) &= P\,g(\symCapitalPsi), \label{eq:A14-P}\\
  g(C\symCapitalPsi) &= C\,g(\symCapitalPsi), \label{eq:A14-C}\\
  g(T\symCapitalPsi) &= T\,g^{-1}(\symCapitalPsi) \label{eq:A14-T}
\end{align}
при том же $(Z_t,Z_{t-1})$, где $g^{-1}$ --- из \Axref{ax:A13}.
Уравнения~\eqref{eq:A14-U1}--\eqref{eq:A14-T} --- точные условия симметрии~$g$.
\end{axiom}

\begin{axiom}[A15 · Без подгоняемых безразмерных констант]
\label{ax:A15}
Безразмерные коэффициенты закона~$g$ --- отношения скоростей, дисперсии, поправки насыщающей фазы,
константы связи на макроуровне --- \emph{не} вводятся как свободные параметры настройки.
Они должны следовать из замыкания условий $A_1$--$A_{16}$ на носителе~$\carrierLattice$,
алфавите регистра и требований~\Axref{ax:A14}.
\end{axiom}

\begin{axiom}[A16 · Фермион / запрет Паули]
\label{ax:A16}
Спин~$1/2$: $R(N_{\mathrm{ring}}/2)\mapsto -z$ (оборот $2\pi$ меняет знак; $4\pi$ --- тождство).
Два \emph{совпадающих} спинора в одном $\PlanckCell$ запрещены (Pauli exclusion на решётке).
Минимальное фазовое действие $\Delta\phi_{\min}=\tfrac12$ задаёт $\ladderAction=\planckHbar/2$.
\end{axiom}

